% !TeX root = ../papers/nc-p3.tex
\section{Walsh基底を持つ四窓の三つの切断と階数12,14,15の認識}
\label{sec:walsh-symmetric-windows}
$E=(\mathbb Z/2)^2=\{0,1,2,3\}$ の加法を $+$、元の低位、高位のビットを
$x_0,x_1$ と書く。四位置に $E$ でラベルされた基底を選び、四次テンソル $W$ が
$i+j+k+t=0$ に台を持つとする。全ラベルへの1の加算は係数を保ち、
2の加算は係数に $(-1)^{(i+k)_0}$ を掛けると仮定する。
以下「Walsh基底を持つ」とはこのような四基底の存在をいう。基底は同型のデータに含めない。
\begin{theorem}\label{thm:walsh-symmetric-windows}
二つの隣接切断がともに階数6のWalsh基底を持つ四次テンソルを取る。
各位置の二次関係として巡回隣接切断の像を用いる。
\begin{enumerate}
\item 反対階数15の全候補は、定理\ref{thm:rank-fifteen-all-ordinary}の境界四点を含む
規格化族にちょうど一致する。従って全候補が標準Koszul左右AS正則で、
指定基底を忘れた同型類と周期は同定理の有理曲面で尽くされる。
\item 反対階数12または14なら、隣接する二切断で非零の文字ブロックがそれぞれ三つあることが
標準Koszul左右AS正則性の必要十分条件である。階数12で正則なら
定理\ref{thm:rank-twelve-quantum-limit}の唯一の量子窓である。
階数14の全正則窓は二本の通常楕円表示とそのWalsh表示の端で覆われる。
\item この対称性を持つ標準窓には反対階数11,13は生じない。
\end{enumerate}
\end{theorem}
\begin{proof}
$b=i+j$, $e=j+k$, $d=b+e$ と置く。対称性から
\[
 W_{i,i+b,i+d,i+e}=(-1)^{i_1d_0}m_{be}
\]
と書ける。$01|23$ の $b$ ブロックは $i,k$ を行列添字とする。
$d$ ごとの符号付き移動 $P_d$ は $(P_d)_{ik}=(-1)^{i_1d_0}[k=i+d]$ である。
$P_1^2=P_2^2=1$, $P_1P_2=-P_2P_1$ から生成代数は $M_2$ となり、
四次元表現はその標準表現の二重である。$s_d=(-1)^{d_0+d_1}$ とすれば
二次元表現の行列式は $\sum_ds_d m_{b,b+d}^2$ である。実際 $P_1,P_2,P_3$ は
対ごとに反可換で、二乗は順に $1,1,-1$ である。$12|30$ の $e$ ブロックを $j,t$ で書くと
符号は $(-1)^{j_1 d_0}(-1)^{(d+e)_1d_0}$ となる。後の因子は各係数の平方を保つので、
列に対して同じ行列式条件を得る。他方 $02|13$ の $d$ ブロックは行の符号を除き
通常の巡回行列で、四固有値は $b\mapsto m_{b,b+d}$ の四つの文字和である。

両隣接切断とも各ブロックが偶数階数なので、最低一つずつ欠ける。
二つ以上欠ける側があれば、全辺で欠けている色を選んだ語は、関係を挿入した
どの語の台にも現れない。各辺の色と最初の文字を独立に選べるので、長さ $n$ で
少なくとも $4\cdot2^{\lfloor(n-1)/2\rfloor}$ 個が独立となり標準成長に反する。
また一方に非零の色が二つ以下なら、各 $d$ の対角列の台は高々二点である。
二点以下の文字和は非零の個数が偶数となるので、階数15では成長を仮定せず三色ずつが必要である。

四位置を $(0,B,B+E',E')$ ずらすと $b,e$ を任意にずらせる。
余分な符号 $(-1)^{i_1(B+E')_0}$ は最初の基底で除く。
両方の欠色を3にしたとき $m$ は左上 $3\times3$ に台を持つ。
全非零ブロックで階数2となる条件は $N_{be}=s_bs_em_{be}^2$ の全行和・列和の消滅である。
特に $N_{01}-N_{10}=N_{12}-N_{21}=N_{20}-N_{02}$ である。
この共通値が零なら $d\ne0$ の各対の係数は符号を除いて一致し、そのブロックは高々階数2、
従って反対階数は高々10である。非零ならこれら三ブロックは可逆で、
$\rho=12+\operatorname{rank} W_{02|13,0}$ となる。
最後の四つの文字和は台が高々三点なので一つだけが非零とはならない
（逆Fourier変換を取れば台が四点になる）。これで第三項も従う。

$\rho=12,14,15$ なら対角の文字和の一つが零である。$m_{be}$ に $b$ の文字を掛ける操作は
位置0,1に同じ対角文字を掛ける基底変換である。零を自明文字に移し、三つの対角成分を
$a+b,-b,-a$ と書く。残りは
\[
 m=\begin{pmatrix}a+b&p&n\\ m'&-b&l\\ -q&r&-a\end{pmatrix},\qquad
 r^2+l^2=2ab,\quad p^2=b^2+r^2,\quad (m')^2=b^2+l^2,\quad
 q^2=a^2+r^2,\quad n^2=a^2+l^2 .
\]
右の五式は全行和・列和と同値であり、$d\ne0$ の各ブロックが可逆なのは $r^2-l^2\ne0$ と同値である。
定理\ref{thm:rank-fifteen-parameter-orbits}のPauli表示の八係数は
\[
 (a,b,L,M,A,B,C,D)=
 \left(a,b,\tfrac{p+m'}2,\tfrac{p-m'}2,\tfrac{n-q}2,\tfrac{n+q}2,
                 \tfrac{l-r}2,\tfrac{l+r}2\right).
\]
四位置のWalsh変換でこれがまさに上の $W$ になることは各添字での文字和で得られる。
\path{computations/verify_walsh_symmetric_windows.py} は成分の恒等式と五式を独立変数で検算する。

階数15なら $ab(a+b)\ne0$ である。$a=1$ に規格化し
$s=q+r$, $t=n+l$, $N'=p+m'$, $H=r^2-l^2$, $K=1+2b$ とする。
$s^{-1}=q-r$, $t^{-1}=n-l$, $H/N'=p-m'$ なので三元は全て単元である。
$r^2+l^2=2b$ より
\[
 K=\tfrac14(s^2+s^{-2}+t^2+t^{-2}),\quad
 H=\tfrac14(s^2+s^{-2}-t^2-t^{-2}),\quad
 (N')^4+(1-K^2)(N')^2+H^2=0 .
\]
$K^2\ne1$, $H\ne0$ で、これが定理\ref{thm:rank-fifteen-rational-quotient}の
規格化表示そのものである。逆は同じWalsh変換から従う。
階数12なら $a=b=0$ であり、八係数の五式に定理\ref{thm:rank-twelve-quantum-limit}を適用できる。

階数14では三つの対角値のちょうど一つが零である。八係数の表示を保つ192変換は
ラベル $1,8,9$ に対応する $b,a,a+b$ のいずれも $b$ の位置へ運べる。
よって $b=0$, $a=1$ として上の同じ復元式を使い、$K=1$, $H\ne0$ を得る。
$H\notin\{1,-1,2,-2\}$ なら倍率は通常であり
（$K+H\ne0,\pm1$, $K-H\ne\pm1$）、定理\ref{thm:rank-fourteen-all-ordinary}による。
残りの二水準 $H^2=1,4$ では平方根の組 $(s,t,N')$ はそれぞれ128,64個である。
三元の逆の式があるので、異なる組は規格化係数として相異なる。

97の上の還元で全ての根を取り、$b=0$ を保つ表示変換の軌道を調べる。
各水準はそれぞれ64個ずつの二軌道、一軌道に分かれる。$H^2=1$ の各軌道は
通常な表示を含む。$H^2=4$ の一代表では全始点で双対経路は $(1,4,6,4,1,0)$、
元の経路は五次まで標準次元、四次中心は五次元で、
四中心商は長さ13の消滅まで所定の次元を持つ。前掲コードは全根とこれらの関係を直接計算する。
この有限集合での軌道計算は標数零にも成立する。実際、$a=1,b=0$ で変換した
$H^2=16(LM)^2$ は $LM=-AB=-CD$ により不変であり、
規格化の分母は四乗根の単元である。$s,t$ は $s^2+s^{-2}=2(1+H)$,
$t^2+t^{-2}=2(1-H)$ の根、$(N')^4=-H^2$ である。
$H=\pm2$ で重なる相反根は既に標数零で重なり、97で他の合流はない。
従って全パラメータ集合の還元は全単射で群作用と可換である。

最後の一代表を持ち上げる際、上の三つの根の方程式は $H$ 上モニックで定数項が単元なので
特殊点は通常な階数14窓の閉包にある（各根の環は有限自由、そのテンソル積も有限自由）。
通常点からの下界と剰余計算の上界が道空間の五次まで、双対経路、全四次中心で一致する。
中心の核は持ち上がり、四中心商の上界と全次数の経路次元の下界から
定理\ref{thm:rank-fourteen-all-ordinary}と同じ自由性・Koszul完全性・左右完全対合が従う。
\end{proof}
